\documentclass[12pt]{article}
\usepackage[utf8]{inputenc}
\usepackage[T1]{fontenc}
\usepackage[margin=1in]{geometry}
\usepackage{setspace}
\usepackage{xcolor}
\usepackage[colorlinks=true, urlcolor=blue!60!black, linkcolor=blue!60!black]{hyperref}
\setlength{\parindent}{0em}
\setlength{\parskip}{0.6em}
\renewcommand{\baselinestretch}{1.0}
\begin{document}
\begin{center}
{\LARGE \textit{New York Times Co. v. United States}, 403 U.S. 713 (1971)\par}
\vspace{0.5em}
{\large Case Brief --- Two-Page Edition\par}
\end{center}

\section{Memory Jogger}
The Pentagon Papers Case: the Government failed to carry the ``heavy burden'' needed to justify a prior restraint on publication (\href{https://supreme.justia.com/cases/federal/us/403/713/}{Justia}).

\section{Detailed Case Facts}
Daniel Ellsberg leaked a classified 47-volume history of U.S. decision-making in Vietnam to the \textit{New York Times} and \textit{Washington Post}. When the \textit{Times} began publishing it in June 1971, the Nixon Administration sued to stop publication on national-security grounds.

\section{Procedural History}
The Second Circuit continued a restraint against the \textit{Times}; the D.C. Circuit refused one against the \textit{Post}. The Supreme Court heard argument on June 26 and ruled on June 30.

\section{Judicial Votes}
6--3 \textit{per curiam}. Black, Douglas, Brennan, Stewart, White, and Marshall each concurred; Burger, Harlan, and Blackmun each dissented.

\section{Holding}
Any system of prior restraints bears a ``heavy presumption against its constitutional validity,'' and the Government did not meet its heavy burden of justification. No rationale commanded a majority.

\section{Analysis of Opinions}
\textbf{Black and Douglas:} the First Amendment forbids prior restraints outright. \textbf{Brennan:} only proof of inevitable, direct, and immediate catastrophe would do. \textbf{Stewart:} the press checks executive power, and no document posed direct, immediate, and irreparable harm. \textbf{White and Marshall:} Congress had not authorized injunctions, though criminal prosecution might follow. \textbf{Dissents:} the Court acted with ``unseemly haste'' and should defer to executive judgments.

\section{Examples: Future Applications}
\textbf{Same-Side:} a leaked historical memo on exaggerated intelligence cannot be enjoined. \textbf{Opposite-Side:} a former analyst bound by a prepublication agreement can be (\textit{Snepp}). \textbf{Fence-Sitter:} a narrow statute that would let courts enjoin publishing the names of covert officers serving abroad.

\section{Critique}
The fractured opinions left the law unsettled, and the threat of Espionage Act prosecution may chill the press more than an injunction would.

\section{Key Quotations}
\begin{quote}
``I am convinced that the Executive is correct with respect to some of the documents involved. But I cannot say that disclosure of any of them will surely result in direct, immediate, and irreparable damage to our Nation or its people. That being so, there can under the First Amendment be but one judicial resolution of the issues before us.''
\end{quote}
\noindent --- Justice Stewart, concurring, 403 U.S. at 730.

\end{document}
